% Clase 16 --- Los dos límites del segmento (§4.3).
% (a) "Imaginemos el más simple de todos, que es el cable subiendo para acá y un
%      cachito después bajando por el otro lado. Entonces las dos contribuciones
%      se me van a compensar."
% (b) "Lo que van a hacer son círculos cerrados... el cable va según el pulgar y
%      los dedos recorren como el campo magnético."
% La geometría es la misma en los dos paneles (un cable vertical y un punto a
% distancia X); lo único que cambia es el régimen: X >> L contra X << L.
% Sentido en (b): corriente hacia arriba (y), a la derecha del cable
% ds x r = y x x = -z, entrante; a la izquierda, saliente; y en el punto de
% adelante (hacia el lector) el campo apunta a la derecha.
\begin{center}
\begin{tikzpicture}[scale=1]

  % =============== (a) lejos: el circuito se cierra ===============
  \begin{scope}
    % circuito cerrado angosto
    \draw[figline, line width=1.6pt] (-0.32,-1.3) rectangle (0.32,1.3);
    \draw[figvecres] (0.32,-0.45) -- (0.32,0.45);
    \draw[figvecres] (-0.32,0.45) -- (-0.32,-0.45);
    \node[figlbl, figred, anchor=west] at (0.4,-0.75) {$I$};
    \node[figblue, scale=1.05] at (0,0) {$\odot$};
    \node[figlblaux, figblue, anchor=south] at (0,1.42) {$\vec\mu = I\vec A$};

    % cota L
    \draw[figvecaux] (-0.75,0.25) -- (-0.75,1.3);
    \draw[figvecaux] (-0.75,-0.25) -- (-0.75,-1.3);
    \node[figlbl, anchor=east] at (-0.85,0) {$L$};

    % distancia al punto lejano
    \draw[figaux] (0.4,0) -- (3.05,0);
    \node[figlbl, anchor=south] at (1.75,0.04) {$X \gg L$};

    % las dos contribuciones, casi iguales y opuestas
    \node[figblue, scale=1.0] at (3.3,0.5) {$\otimes$};
    \node[figblue, scale=1.0] at (3.3,-0.5) {$\odot$};
    \node[figlblaux, anchor=west, align=left] at (3.52,0.5) {del que sube};
    \node[figlblaux, anchor=west, align=left] at (3.52,-0.5) {del que baja};
    \node[figlbl, anchor=north, align=center] at (2.6,-1.05)
      {casi se cancelan};

    \node[figlbl, anchor=north, align=center] at (2.0,-1.95)
      {un tramo solo daría $B\sim\dfrac{1}{X^2}$,\\[0.25em]
       el circuito cerrado deja $B\sim\dfrac{1}{X^3}$};

    \node[figlbl, anchor=north, align=center] at (2.0,-3.95)
      {(a) lejos: es un \textbf{dipolo} magnético};
  \end{scope}

  % =============== (b) cerca: hilo infinito ===============
  \begin{scope}[xshift=8.3cm]
    % el cable, que se prolonga fuera del dibujo
    \draw[figline, line width=1.8pt] (0,-2.15) -- (0,2.15);
    \draw[figline, line width=1.8pt, dotted] (0,2.15) -- (0,2.6);
    \draw[figline, line width=1.8pt, dotted] (0,-2.15) -- (0,-2.6);
    \draw[figvecres] (0,1.45) -- (0,2.05);
    \node[figlbl, figred, anchor=west] at (0.1,1.75) {$I$};

    % líneas de campo: círculos cerrados alrededor del cable
    \foreach \yy in {-1.05,0,1.05} {
      \draw[figaux, figblue] (1.5,\yy) arc (0:180:1.5 and 0.4);
      \draw[figline, line width=0.7pt] (-1.5,\yy) arc (180:360:1.5 and 0.4);
      \draw[figvec, line width=0.7pt] (-0.32,{\yy-0.4}) -- (0.36,{\yy-0.4});
    }

    % el campo a un lado y al otro
    \node[figblue, scale=1.05] at (2.15,0.42) {$\otimes$};
    \node[figblue, scale=1.05] at (-2.15,0.42) {$\odot$};

    % cota X
    \draw[figaux] (1.5,-1.45) -- (1.5,-1.95);
    \draw[figvecaux] (0.65,-1.85) -- (0,-1.85);
    \draw[figvecaux] (0.85,-1.85) -- (1.5,-1.85);
    \node[figlbl, anchor=south] at (0.75,-1.8) {$X$};

    \node[figlbl, anchor=north] at (0,-2.65)
      {$B = \dfrac{\mu_0 I}{2\pi X}$};

    \node[figlbl, anchor=north, align=center] at (0,-3.95)
      {(b) cerca: es un \textbf{hilo infinito}};
  \end{scope}

\end{tikzpicture}\\[0.5em]
{\small\itshape El resultado exacto del segmento tiene dos regímenes. Visto de
lejos ($X\gg L$) da $B\simeq\mu_0 IL/4\pi X^2$, y ahí conviene desconfiar: ese
$1/X^2$ es el de un tramo que, solo, no puede existir. Cerrando el circuito
---aunque sea con el más simple, un tramo que sube y otro que baja--- las dos
contribuciones llegan al punto lejano casi iguales y de sentido opuesto, y lo
que sobrevive decae como $1/X^3$. La razón de fondo es que no hay monopolos
magnéticos: no existe el análogo de la carga aislada, así que lo primero que se
ve de lejos es un \textbf{dipolo} de momento $\vec\mu = I\vec A$, exactamente
como el campo de un dipolo eléctrico y a diferencia del $1/r^2$ de una carga
puntual. En el otro extremo ($X\ll L$) el segmento se comporta como un hilo
infinito y queda $B=\mu_0 I/2\pi X$, con las líneas de campo formando círculos
cerrados alrededor del cable, tangentes en todo punto y orientados por la regla
de la mano derecha con el pulgar según la corriente. Ese caso es el punto de
partida de la próxima clase.}
\end{center}
